\section{Higher-rank symplectic cores in every odd field at least seven}
\label{sec:higher-odd-symplectic}

\paragraph{Status and precise scope.}
This complete new derivation has passed whole author, root and
independent peer proof reviews, with separately executed exact certificates. It uses the proved elementary carrier and degree
lemmas D-HOS-01--02. The only character-ratio input is the published
unipotent estimate stated below. No full character table, classification
of small character degrees, or Weil-character formula is a premise.
The final theorem permits mixed ranks and fields and does not require
transitivity on the factors, essential selfduality, or a polarization
determinant condition. The odd fields $3$ and $5$ are outside its scope.

For a finite-order complex operator $T$ of degree $a$, put
\[
 \nu(T)=a-\max_\lambda\dim\ker(T-\lambda I).
\]
Scalar multiplication does not change this quantity. Ordinary moving
rank is at least $\nu(T)$, and $\nu(T)\ge\nu(T^j)$ for every integer
$j$. For a unitary operator, the triangle inequality at its largest
eigenspace gives
\[
                 |\operatorname{Tr}T|\ge a-2\nu(T).       \tag{HOS.13}
\]

The exact published input is
\cite[simply connected conventions, p.~2;
Theorem~2.4(i), pp.~3--4]{HES:Liebeck2017}, the published restatement
of Gluck's unipotent estimate. For $q>4$, a nontrivial ordinary
irreducible character $\chi$ of $\operatorname{Sp}_{2n}(q)$ satisfies
\[
 |\chi(u)|\le c_q\chi(1),\qquad
 c_q=\frac1{\sqrt q-1},\qquad
                 1\ne u\text{ unipotent}.                 \tag{HOS.14}
\]
The subtraction is outside the radical. In what follows $q\ge7$,
so $0<c_q<1$. Define
\[
 \delta_q=\frac12(1-q^{-1})(1-q^{-3})(1-c_q),\qquad
 \delta_* =\frac{9234}{55223}>\frac16.                     \tag{HOS.15}
\]

\begin{lemma}[D-HOS-03: every inner operator]
\label{D-HOS-03}
Let $q\ge7$ be odd, $n\ge4$, and let $\rho$ be a nontrivial
ordinary irreducible representation of $S=\operatorname{Sp}_{2n}(q)$
of degree $a$. For every noncentral $x\in S$,
\[
                         \nu(\rho(x))\ge\delta_q a.        \tag{HOS.16}
\]
\end{lemma}

\begin{proof}
The projective kernel of $\rho$ is exactly $Z(S)$: its image in the
simple quotient is normal, and if this image were the whole quotient,
the perfect group would act by scalars and hence trivially. Take a
power $h$ of $x$ whose image in $S/Z(S)$ has prime order $\ell$.
Its support is no greater than that of $x$. If $\ell=p$, where
$q=p^f$, multiplication by a central sign makes $h$ a nonidentity
unipotent. Equations (HOS.13)--(HOS.14) give support at least
$(1-c_q)a/2\ge\delta_q a$.

Suppose $\ell\ne p$. The element $h$ is semisimple. If $\ell$
is odd, a central sign makes its actual order $\ell$. If $\ell=2$,
its square can be either $I$ or $-I$. We retain this latter
possibility throughout the argument.

Let $E$ be the natural symplectic space with alternating form
$B$. First suppose $h$ has a nonzero invariant totally isotropic
subspace $R$. Choose a basis $r_1,\ldots,r_d$ of $R$. For a
symmetric matrix $C=(c_{ij})$ set
\[
 N_C(v)=\sum_{i,j}c_{ij}B(r_j,v)r_i,\qquad u_C=I+N_C.
                                                               \tag{HOS.17}
\]
Isotropy gives $N_CN_D=0$, and symmetry gives
$B(N_Cv,w)+B(v,N_Cw)=0$. Thus the $u_C$ form an additive
unipotent subgroup of $S$. The map $C\mapsto N_C$ is injective,
because the pairing functionals of the independent $r_j$ are
independent. If $A=h|_R$, conjugation is $C\mapsto ACA^t$.
The action has order dividing $\ell$, including when $h^2=-I$.

If this action is nontrivial, its moving subspace over
$\mathbf F_q$ has no fixed vectors: an operator of prime order
coprime to the characteristic is semisimple, and its moving and
fixed subspaces are complementary. It gives a subgroup $U$ of
order $Q\ge q$. Every nonzero dual weight has an orbit of length
$\ell$. Averaging (HOS.14) bounds its trivial weight dimension by
$a(c_q+(1-c_q)/Q)$. On each other orbit a cyclic block operator
has every eigenspace at most $1/\ell$ of that orbit's dimension;
this conclusion does not require its global $\ell$th power to be
scalar. Therefore
\[
 \nu(\rho(h))\ge(1-1/\ell)(1-Q^{-1})(1-c_q)a
       \ge\tfrac12(1-q^{-1})(1-c_q)a\ge\delta_q a.         \tag{HOS.18}
\]

It remains in this isotropic case that $ACA^t=C$ for all symmetric
$C$. Since $A$ is semisimple, its eigenvalue squares and all their
pairwise products equal one. Hence $A=\lambda I$ with
$\lambda=\pm1$. In particular $h^2=-I$ is impossible in this
case. Choose a nonzero $e\in R$. There is a paired vector $f$
with $B(e,f)=1$ and $hf=\lambda f$: project a paired vector to
the $\lambda$ eigenspace by the coprime-order averaging projector.
Explicitly, if $o$ is the actual order of $h$, use
$o^{-1}\sum_{j=0}^{o-1}\lambda^j h^j f$.
The plane $\langle e,f\rangle$ is nondegenerate, and its orthogonal
complement $W$ is $h$-stable. Put $g=\lambda h|_W$. This is
symplectic of order dividing $\ell$ and is nonidentity, since
otherwise $h=\lambda I$ on all of $E$. Its moving subspace
$W_0=\operatorname{im}(g-I)$ is nondegenerate, has even dimension
$2m\ge2$, and $g$ has no fixed vectors on it.

The line-parabolic radical contains the following subgroup, with
$v\in W_0$ and $z\in\mathbf F_q$:
\[
 e\mapsto e,\qquad w\mapsto w+B(v,w)e,\qquad
                 f\mapsto f+v+2ze.                         \tag{HOS.19}
\]
Its multiplication is the Heisenberg law
$(v,z)(v',z')=(v+v',z+z'+B(v,v')/2)$.
Conjugation by $h$ is $(v,z)\mapsto(gv,z)$, since
$\lambda^2=1$. Every nonidentity element is unipotent.
D-HOS-01, applied to the full restriction of $\rho$, gives
\[
 \nu(\rho(h))\ge
 (1-1/\ell)(1-q^{-m})(1-q^{-(2m+1)})(1-c_q)a
                         \ge\delta_q a.                    \tag{HOS.20}
\]
The lemma includes both linear-weight orbits and all nontrivial
central-character carriers with arbitrary multiplicity operators.
Thus no assumption about the restriction being irreducible or
Weil is needed here.

Finally suppose that $h$ has no invariant nonzero totally isotropic
subspace. In its semisimple primary decomposition, a pair of distinct
reciprocal factors provides such a subspace. A plus- or minus-one
primary summand does too. A self-reciprocal irreducible factor of
degree $2d$ and multiplicity at least two also provides one: the
associated form on the multiplicity space is Hermitian over a finite
quadratic extension, and a Hermitian space of dimension at least two
is isotropic. More explicitly, put $K=\mathbf F_q[h]$ on that
primary space. Its adjoint involution sends $h$ to $h^{-1}$;
its fixed field is $\mathbf F_{q^d}$. Defining a $K$-valued
form by $\operatorname{Tr}_{K/\mathbf F_q}(a b(v,w))=B(av,w)$
for every $a\in K$ gives a nondegenerate skew-Hermitian form,
by the trace-pairing argument written out in (HOS.23) below.
Multiplication by a nonzero scalar negated by the involution
makes it Hermitian. To see its isotropy directly, diagonalize two coordinates
and solve their norm equation using surjectivity of the finite-field
norm. The resulting extension-field line is isotropic for the
original alternating form as well.

Consequently all primary factors are distinct self-reciprocal
irreducibles of degrees $2d_i$, each with multiplicity one, where
$\sum_i d_i=n$. Their commuting natural operators are field scalars,
and preservation of the form imposes the norm-one condition.
It follows that
\[
 |C_S(h)|=\prod_i(q^{d_i}+1)\le(q+1)^n.
\]
Column orthogonality and D-HOS-02 now give
\[
 \frac{|\chi(h)|}{a}
 \le\frac{2(q+1)^{n/2}}{(q-1)q^{n-1}}
 \le\frac{2(q+1)^2}{(q-1)q^3}<\frac1q<c_q.                 \tag{HOS.21}
\]
The middle inequality follows by increasing $n$ from four, each
increment multiplying the expression by $\sqrt{q+1}/q<1$.
For the strict bound use $q+1\le8q/7$ and $q-1\ge6$, so
$2(q+1)^2/((q-1)q^2)\le128/294<1$.
Equations (HOS.13) and (HOS.21) give at least
$(1-c_q)a/2\ge\delta_q a$. All cases have been exhausted.
\end{proof}

\begin{lemma}[D-HOS-04: diagonal projective involutions]
\label{D-HOS-04}
In the setting of D-HOS-03, let $T$ implement a diagonal outer
automorphism whose square is the identity automorphism. Then
\[
                    \nu(T)\ge\tfrac12(1-q^{-1})(1-c_q)a.
                                                               \tag{HOS.22}
\]
\end{lemma}

\begin{proof}
Take a natural lift $h\in\operatorname{GSp}_{2n}(q)$ with
nonsquare multiplier $\mu$. The identity square of the automorphism
implies $h^2=cI$ for some $c\in\mathbf F_q^*$, since the natural
centralizer of the symplectic group is scalar. Thus $\mu=\pm c$.
We construct a nonzero $h$-stable totally isotropic subspace $R$.

If $c$ is a square, $h$ has eigenspaces for the two square roots
of $c$. Since $\mu$ is nonsquare, $\mu=-c$. The form vanishes
on either eigenspace, because its multiplier there would otherwise
be $c$ instead of $-c$. Choose a nonzero eigenline as $R$.

If $c$ is not a square, put $K=\mathbf F_q[u]/(u^2-c)$ and
let $u$ act as $h$ on $E$. The natural space has dimension $n$
over $K$. For $v,w\in E$ define the unique $K$-valued form
$b(v,w)$ by
\[
       \operatorname{Tr}_{K/\mathbf F_q}(a b(v,w))=B(av,w)
                    \quad\hbox{for every }a\in K.          \tag{HOS.23}
\]
Nondegeneracy of the trace pairing gives uniqueness and existence.
The adjoint of $h$ for $B$ is $(\mu/c)h$.
If $\mu=c$, the adjoint involution on $K$ is the identity.
Equation (HOS.23) is then $K$-bilinear and alternating:
$B(av,v)=B(v,av)=-B(av,v)$, so it vanishes for every $a$.
It is nondegenerate, and any $K$-line is isotropic. This case
automatically has $n$ even.

If $\mu=-c$, the adjoint involution on $K$ is its nontrivial
Galois involution. Equation (HOS.23) is linear in its first
variable, conjugate-linear in its second, and satisfies
$b(w,v)=-\overline{b(v,w)}$. Multiplying by any nonzero
trace-zero scalar makes it Hermitian. It is nondegenerate;
as $n\ge4$, its finite-field norm equation supplies an isotropic
$K$-line. In either case such a line, viewed over $\mathbf F_q$,
is an $h$-stable isotropic plane $R$ for $B$.

Use the shear subgroup (HOS.17) for this $R$. Conjugation is
now the normalized action
\[
                         C\longmapsto\mu^{-1}ACA^t,
                              \qquad A=h|_R.                \tag{HOS.24}
\]
This follows directly by moving $h$ through the pairing in
(HOS.17), which introduces the factor $\mu^{-1}$.
The action has square one. It is not the identity: otherwise,
over an algebraic closure, every eigenvalue square of the
semisimple $A$ and every pairwise product would be $\mu$.
All eigenvalues would coincide, so $A=\alpha I$ with
$\alpha\in\mathbf F_q$ and $\alpha^2=\mu$, contradicting
the nonsquare multiplier. Its minus-one eigenspace gives an
additive unipotent subgroup $U$ of order at least $q$, on which
$T$ exchanges all nonzero dual weights in pairs.
The averaging and cyclic-block calculation in (HOS.18), with
$\ell=2$, proves (HOS.22). The proof did not require $T^2$
to be the identity as an operator, only the indicated action
on the core and its shear subgroup.
\end{proof}

\begin{theorem}[D-HOS-05: every compatible normalizer]
\label{D-HOS-05}
For $q\ge7$ odd and $n\ge4$, every finite-order nonscalar
operator normalizing the image of a nontrivial ordinary irreducible
representation of $\operatorname{Sp}_{2n}(q)$ of degree $a$ satisfies
\[
                              \nu(T)\ge\delta_* a.          \tag{HOS.25}
\]
All scalar, diagonal and field possibilities are included.
The same conclusion holds for every quasisimple central quotient
with this simple quotient.
\end{theorem}

\begin{proof}
There are no exceptional multipliers in this range, by
\cite[Chapter~5, complete Table~1 and preceding paragraph,
p.~35]{BI:HHM}. Thus the universal perfect cover is the symplectic
group. An automorphism of a represented central quotient lifts
uniquely to that cover; an automorphism trivial on the simple
quotient is trivial, since its difference from the identity is a
homomorphism from a perfect group to its center. The exact
outer group is $C_2\times C_f$ for $q=p^f$, generated by diagonal
and field classes, from \cite[\S2.2.2, Table~2.1,
author p.~14]{HES:LiXia2022}. These descriptions therefore
apply to every normalizing operator under consideration.

The function $\delta_q$ is increasing for $q\ge7$, since each
of its three positive factors is increasing. Also
$\sqrt7>37/14$, because $1372>1369$. Hence
$c_7<14/23$ and
\[
 \delta_7>\tfrac12\frac67\frac{342}{343}\frac9{23}
                  =\frac{9234}{55223}=\delta_* .            \tag{HOS.26}
\]
Thus every noncentral inner operator already satisfies (HOS.25).

If $q$ is prime, there is no field component. For a diagonal
outer operator, if the square of its induced automorphism is a
nontrivial inner automorphism, $T^2$ is scalar times a noncentral
inner operator, and D-HOS-03 applies. If the square is the
identity, D-HOS-04 gives an even stronger estimate.

For every $q$, a general normalizer has support at least
$\delta_q a/2$. Indeed, a nonscalar $T$ fails to commute with
some represented core element by Schur's lemma. A noncentral
commutator has moving rank at least $\delta_q a$ by D-HOS-03;
a nonidentity scalar commutator has moving rank $a$.
After scalar-normalizing $T$ at its largest eigenspace,
\[
             \operatorname{rank}([T,g]-I)\le2\nu(T),
\]
which proves the claim in both cases. Consequently all $q\ge25$
are covered, because
\[
 \frac{\delta_q}{2}\ge\frac{\delta_{25}}2
 =\frac9{50}\frac{15624}{15625}>\delta_* .                  \tag{HOS.27}
\]
Every nonprime odd prime power at least seven is either nine
or at least twenty-five. It remains only to handle $q=9$.

The outer group at $q=9$ is $C_2\times C_2$. If the square of
the induced automorphism is nontrivial inner, use $T^2$ as
above. A pure diagonal involution is covered by D-HOS-04.
Consider a field-containing automorphism whose square is the
identity. It has a natural lift $h\sigma$, where
$h\in\operatorname{GSp}_{2n}(9)$ and $\sigma$ is coordinatewise
cubing. Its square is scalar, say
\[
                         h\sigma(h)=cI.
\]
Associativity gives $c=\sigma(c)$, hence $c\in\mathbf F_3^*$.
Writing $\mu$ for the multiplier of $h$ gives
$\mu\sigma(\mu)=c^2=1$. In particular $\mu$ is a square
in $\mathbf F_9^*$; a mixed diagonal-field class cannot have
identity square. We nevertheless allow an arbitrary such lift
in the following descent, so no representation is assumed
invariant under a separate diagonal conjugation.

Choose $z\in\mathbf F_9^*$ with $z\sigma(z)=c^{-1}$ and set
$s=zh\sigma$. Then $s^2=I$. Its multiplier coefficient
$\mu'=z^2\mu$ has norm one. Choose $d\in\mathbf F_9^*$ with
$\sigma(d)/d=\mu'$. Both choices follow from the cyclic
multiplicative group: the norm has image $\mathbf F_3^*$,
and $d\mapsto\sigma(d)/d=d^2$ has image the norm-one subgroup.
For $B'=dB$ one has
\[
                         B'(sv,sw)=\sigma(B'(v,w)).          \tag{HOS.28}
\]
The fixed space $E_0=E^s$ is $2n$-dimensional over $\mathbf F_3$
and spans $E$ over $\mathbf F_9$. Explicitly, choose
$\alpha\in\mathbf F_9$ with $\alpha^2=-1$, so
$\sigma(\alpha)=-\alpha$. Every vector has the decomposition
\[
             v=\frac{v+sv}{2}
                 +\alpha\frac{v-sv}{2\alpha},              \tag{HOS.29}
\]
whose two displayed coefficients belong to $E_0$.
Moreover $E_0\cap\alpha E_0=0$, which proves the dimension
and independence assertion. The restriction of $B'$ to $E_0$
is an alternating form over $\mathbf F_3$ and is nondegenerate,
since $E_0$ spans. A symplectic basis supplies a matrix
$P\in\operatorname{GSp}_{2n}(9)$ with
$P^{-1}sP=\sigma$. Replace the representation by
$x\mapsto\rho(PxP^{-1})$. Its intertwiner implements pure
$\sigma$, and it has the same degree and the same unipotent
estimate. This relabeling is permissible even if $P$ has
nontrivial diagonal class.

In the standard $\operatorname{GL}_n(9)$ Levi choose the
one-row subgroup
$U=\{\left(\begin{smallmatrix}1&u\\0&I_{n-1}\end{smallmatrix}\right)
:u\in\mathbf F_9^{\,n-1}\}$.
The block subgroup $\operatorname{diag}(1,\operatorname{GL}_{n-1}(9))$
normalizes $U$ and is transitive on its nonzero dual weights,
so each has the same multiplicity $m$. Put $Q=9^{n-1}$ and
$M=3^{n-1}$.
Since $c_9=1/2$, averaging (HOS.14) gives $m\ge a/(2Q)$.
Frobenius fixes exactly $M$ of the $Q$ dual weights, under
the nondegenerate trace-pairing identification of that dual.
All the remaining weights occur in pairs. Therefore
\[
 \nu(T)\ge\frac{Q-M}{2}m
      \ge\frac14(1-3^{-(n-1)})a
      \ge\frac{13}{54}a>\delta_*a.                          \tag{HOS.30}
\]
This proves the last field case. The same lifting argument
from the universal perfect cover proves the central-quotient
assertion. Every compatible finite-order normalizer is covered.
\end{proof}

\begin{theorem}[D-HOS-06: the higher odd symplectic source exclusion]
\label{D-HOS-06}
Suppose an inducing block in the signature-$(3,3)$ sixfold
curve-cover problem has an irreducibly acting normal central product
of quasisimple groups, each with simple quotient
$\operatorname{PSp}_{2n_i}(q_i)$, where $n_i\ge4$ and
$q_i\ge7$ is odd. Then the actual period image has dimension
at most eight. This allows different ranks and fields, arbitrary
compatible scalar and outer operators, and all factor permutations.
It does not require equality of the inducing and original character
fields, selfduality of any representation, or a determinant condition.
\end{theorem}

\begin{proof}
The core representation is a tensor product of nontrivial irreducible
factor modules. Let $e$ be its total degree. For a factor-preserving
nonscalar operator, its normalized eigenspace multiplicities are
convolutions. Their maximum is bounded by the maximum for any
nonscalar factor, so D-HOS-05 gives support at least $\delta_*e$.

For clarity, the same estimate holds when factors are permuted.
On a nontrivial factor cycle of length $r$ and factor degree $a$,
conjugate the intertwiners to identities except for one closing
twist and diagonalize that twist. The resulting operator is a
weighted cyclic permutation of the $a^r$ basis words. There are
$a$ constant words, and all other word orbits have length at least
two. Each orbit contributes at most one to any chosen eigenspace.
Thus the largest eigenspace is at most
$a+(a^r-a)/2\le3a^r/4$. Tensoring with the other cycles cannot
increase the normalized maximum. The resulting support $e/4$
is stronger than $\delta_*e$. Nonidentity scalar inertias move
all $e$ dimensions.

D-HOS-02 gives for every factor
\[
             a_i\ge\frac{q_i-1}{2}q_i^{n_i-1}
                       \ge3\cdot7^3=1029,
\]
so $e\ge1029$. Apply the actual intermediate-cover factorization:
if its base has genus $h$ and $R$ genuine inertia points, then
\[
              6=2e(h-1)+D,\qquad
                  D=\sum_{j=1}^{R}\operatorname{rank}(g_j-I).
                                                               \tag{HOS.31}
\]
This is the already established polarized multiplicity variation
factorization, valid also for a larger inducing character field.
Genus at least two is impossible because $D\ge0$ and $e>3$.
In genus one, $D=6$, hence $R\le6$ and the actual pointed-base
Hurwitz variation has dimension at most six. A period image of
dimension nine would consequently force genus zero and
$R\ge12$, with $D=2e+6$.

Every genuine inertia has moving rank at least $\delta_*e$.
It follows that
\[
       2e+6\ge12\delta_*e,\qquad
       e\le\frac6{12\delta_*-2}
           =\frac{165669}{181}<916,                          \tag{HOS.32}
\]
contradicting $e\ge1029$. This is a necessary actual-period
bound and does not infer dominance from parameter counts.
\end{proof}

\paragraph{Remaining fields and source boundary.}
The argument makes no assertion for higher-rank odd symplectic factors
over fields $3$ or $5$. The low-rank $C_2$ and $C_3$ results retain
their separately checked scopes. The published inputs here are precisely
the ordinary unipotent ratio, the ordinary automorphism description,
and the complete exceptional-multiplier list. All carrier, isotropic,
quadratic-form, semilinear-descent, tensor, and numerical arguments
are supplied in full above or in D-HOS-01--02. No missing semisimple
ratio theorem, unrefereed result, or full generic character table is used.
